\section{Multi-Layer Analysis and Optimal Depth}
\label{sec:multi}

We now extend the analysis to multiple layers, derive the optimal organizational depth, and characterize the optimal budget allocation.

\subsection{The Endogenous Transmission Factor}
\label{subsec:transmission_factor}

Before analyzing information transmission across layers, we formally define the endogenous transmission factor and relate it to the single-layer distortion measure.

Recall from Definition~\ref{def:single_layer_distortion} that the single-layer distortion measure is $\tilde{\delta}_\ell = 1 - \Delta_\ell$. In the multi-layer context, it is more convenient to work directly with the ratio of achieved posterior precision to first-best precision:

\begin{definition}[Endogenous Transmission Factor]\label{def:transmission_factor}
The endogenous transmission factor from layer $\ell$ to layer $\ell+1$ is
\begin{equation}\label{eq:transmission_factor}
\rho_\ell \equiv \frac{\tau^*_\ell}{\tau^{\text{FB}}_\ell} = \frac{\tau^{\text{prior}}_\ell + \sum_{k=1}^{K_\ell} g_\ell(\alpha^*_{\ell,k}) \tau^0_{\ell,k}}{\tau^{\text{prior}}_\ell + \sum_{k=1}^{K_\ell} \tau^0_{\ell,k}} \in [0,1],
\end{equation}
where $\tau^{\text{FB}}_\ell$ is the first-best posterior precision and $\tau^*_\ell$ is the achieved posterior precision under the optimal attention allocation.
\end{definition}

The transmission factor $\rho_\ell$ measures the fraction of first-best precision that survives layer $\ell$. When $\rho_\ell = 1$, all signals receive sufficient attention and no precision is lost. When $\rho_\ell < 1$, the attention budget forces the layer to ignore or under-process some signals, creating a precision leakage.

\begin{remark}[From Exogenous to Endogenous]\label{rem:exog_to_endog}
In previous work, the transmission factor was an exogenous parameter $\tau \in [0,1]$ that the modeler specified. In our framework, $\rho_\ell$ is endogenously determined by the attention-allocation problem: it depends on the budget $A_\ell$, the number of signals $K_\ell$, their precisions $\{\tau^0_{\ell,k}\}$, and the attention technology $g_\ell(\cdot)$. Importantly, $\rho_\ell$ is derived \emph{relative to the assumed attention technology $g_\ell(\cdot)$}---changing the functional form of $g_\ell$ would alter the numerical value of $\rho_\ell$ even holding all primitives fixed.
\end{remark}

\begin{remark}[Technology Dependence and Qualitative Robustness]\label{rem:rho_technology}
The numerical value of $\rho_\ell$ depends on the specific attention technology $g_\ell(\cdot)$. Under the power-function specification $g_\ell(\alpha) = \alpha^\gamma$, a larger context window (higher $A_\ell$) raises $\rho_\ell$; under an alternative technology, the quantitative response may differ. However, the qualitative conclusions of our analysis---that optimal depth is finite and that budget allocation is front-loaded---do not depend on the specific functional form of $g_\ell(\cdot)$. All that is required is that $g_\ell$ be increasing and concave with $g'_\ell(0) > 0$.
\end{remark}

\subsection{Organizational Design: Optimal Budget Allocation}
\label{subsec:budget_allocation}

Before analyzing the finite-depth result, we first characterize the optimal budget allocation across layers. This is the model's most empirically testable prediction.

Suppose the principal faces a total attention budget $B > 0$ that can be allocated across layers. For simplicity, consider a uniform technology profile where $g_\ell(\cdot) = g(\cdot)$, $c_\ell(\cdot) = c(\cdot)$, and $K_\ell = K$ for all $\ell$. The principal's problem is:
\begin{equation}\label{eq:budget_allocation}
\max_{L, \{A_\ell\}} V_L(\tau^*_L) \quad \text{s.t.} \quad \sum_{\ell=0}^{L} A_\ell \leq B, \quad A_\ell \geq 0.
\end{equation}

We compare three canonical budget-allocation strategies:

\textbf{Strategy 1: Uniform allocation.} Set $A_\ell = B/(L+1)$ for all $\ell$. This is the simplest design: every layer receives the same context-window capacity.

\textbf{Strategy 2: Front-loaded allocation.} Set $A_0 > A_1 > \cdots > A_L$. The bottom layers receive larger budgets because they process raw signals with the highest intrinsic precision. The intuition is that precision loss is most costly early in the chain: a unit of precision lost at layer 0 is compounded by all subsequent layers.

\textbf{Strategy 3: Back-loaded allocation.} Set $A_0 < A_1 < \cdots < A_L$. The top layers receive larger budgets because they make the final decision and their precision directly enters the payoff.

\begin{proposition}[Optimal Budget Allocation]\label{prop:front_loading}
Under the uniform-technology assumption, the optimal budget allocation is front-loaded: $A^*_0 \geq A^*_1 \geq \cdots \geq A^*_L$.
\end{proposition}

\begin{proof}
Consider two adjacent layers $\ell$ and $\ell+1$ with budgets $A_\ell$ and $A_{\ell+1}$. Define an alternative allocation that transfers $\varepsilon > 0$ from layer $\ell+1$ to layer $\ell$: $A'_\ell = A_\ell + \varepsilon$, $A'_{\ell+1} = A_{\ell+1} - \varepsilon$. We show that if $A_\ell = A_{\ell+1}$, then $\tau^*_L(A'_\ell, A'_{\ell+1}) > \tau^*_L(A_\ell, A_{\ell+1})$ for small $\varepsilon$.

The transmission factor at layer $\ell$ increases: $\rho_\ell(A'_\ell) > \rho_\ell(A_\ell)$ because $\partial \rho_\ell / \partial A_\ell > 0$ (Proposition~\ref{prop:distortion_comp_statics}). This raises the inherited precision at all downstream layers $j > \ell$.

The transmission factor at layer $\ell+1$ decreases: $\rho_{\ell+1}(A'_{\ell+1}) < \rho_{\ell+1}(A_{\ell+1})$. However, the inherited precision at layer $\ell+1$ is $\tau^{\text{prior}}_{\ell+1} = \rho_\ell \cdot \tau^*_\ell$, which increases because $\rho_\ell$ increases. For small $\varepsilon$, the increase in $\rho_\ell$ dominates the decrease in $\rho_{\ell+1}$ because the compounding effect amplifies the gain at layer $\ell$ across all downstream layers.

Repeating the exchange argument for all adjacent pairs from $\ell = 0$ to $\ell = L-1$, any non-front-loaded allocation can be improved by shifting budget toward earlier layers. The only allocation immune to such improvements is the front-loaded one.
\end{proof}

\vspace{1em}
\noindent\textbf{Practical Implications.}
In LLM pipeline design, Proposition~\ref{prop:front_loading} implies that the largest model (largest context window) should be deployed at the earliest stage---the layer that processes the raw user query and initial retrieved documents. Smaller, faster models can be used at intermediate layers.

This prediction aligns with the architecture of several prominent production multi-agent systems. \citet{openai2025deepresearch} deploys a large reasoning model (o3, 200K context window) for the initial query-understanding and web-retrieval phase, while using lighter models for intermediate synthesis steps. \textit{Manus} \citep{manus2025} assigns its largest model to the initial task-decomposition layer, where raw user intent must be parsed and mapped to a tool-execution plan. In open-source tool-use frameworks, \textit{OpenHands} \citep{openhands2024} and \textit{Devin} \citep{cognition2024devin} both place a frontier model at the entry point of the pipeline (code comprehension and planning), while delegating implementation steps to smaller specialized agents. Our model provides a theoretical justification for this engineering practice: the returns to attention capacity are highest at the bottom of the chain because precision lost there is amplified by every subsequent layer.

\begin{remark}[Front-loading as an Empirical Signature]\label{rem:front_loading_empirical}
The front-loading result is the model's most directly testable prediction. It does not require measuring the preservation parameters $(\underline{\eta}, \bar{\eta}, a)$ or the transmission factor $\rho_\ell$, which are difficult to observe in practice. Instead, it predicts a monotonic relationship between layer index and budget allocation that can be evaluated using only observable architectural choices: the model sizes or context-window lengths assigned to each layer of a multi-agent pipeline. We return to this prediction in Section~\ref{sec:empirics}, where we discuss how to bring it to data.
\end{remark}

\subsection{Information Transmission Across Layers}
\label{subsec:info_transmission}

Having characterized the optimal budget allocation, we now study how information propagates through the chain. This analysis reveals the dynamic forces that ultimately bound the optimal depth.

In a chain with $L$ layers, the posterior precision at layer $\ell$ depends on the attention allocation at layer $\ell-1$ through the endogenous transmission factor $\rho_{\ell-1}$. The prior precision at layer $\ell$ is the achieved posterior precision at layer $\ell-1$:
\begin{equation}\label{eq:prior_update}
\tau^{\text{prior}}_\ell = \rho_{\ell-1} \cdot \tau^*_{\ell-1}(A_{\ell-1}), \quad \ell = 1, \ldots, L.
\end{equation}
With $\tau^{\text{prior}}_0 \equiv \tau^0_0$ (the raw signal precision at the bottom layer), the posterior precision at layer $\ell$ follows the recursive Bayesian-update structure:
\begin{equation}\label{eq:recursive_precision}
\tau^*_\ell = \tau^{\text{prior}}_\ell + \sum_{k=1}^{K_\ell} g_\ell(\alpha^*_{\ell,k}) \tau^0_{\ell,k}, \quad \text{with } \tau^{\text{prior}}_\ell = \rho_{\ell-1} \cdot \tau^*_{\ell-1}.
\end{equation}

Solving the recursion yields the cumulative precision at the top layer:
\begin{equation}\label{eq:top_precision}
\tau^*_L = \sum_{\ell=1}^{L} \sum_{k=1}^{K_\ell} g_\ell(\alpha^*_{\ell,k}) \tau^0_{\ell,k} \cdot \prod_{j=\ell}^{L-1} \rho_j + \tau^0_0 \cdot \prod_{j=0}^{L-1} \rho_j.
\end{equation}
Each layer $\ell$ contributes its local signal precision $g_\ell(\alpha^*_{\ell,k}) \tau^0_{\ell,k}$, weighted by the subsequent transmission losses $\prod_{j=\ell}^{L-1} \rho_j$.

\vspace{1em}
\noindent\textbf{Precision Dynamics under Information Depreciation.}
We now characterize the dynamics of the precision path $\{\tau^*_\ell\}$ under the information depreciation model introduced in Assumptions~\ref{ass:state_signals} and~\ref{ass:endogenous_eta}. Recall that in this framework, the raw signal quality at layer $\ell$ is $\tau^0_{\ell,k} = \tau^0_k \cdot \eta_\ell^{\,\ell}$, where $\eta_\ell \in (\underline{\eta}, \bar{\eta})$ is the endogenous per-layer signal preservation rate at layer $\ell$, determined by the attention budget $A_\ell$ through equation~\eqref{eq:endogenous_eta}.

To build intuition, consider first what happens when $\eta_\ell = 1$ for all $\ell$ (no signal degradation). Under uniform technology and uniform attention budgets ($g_\ell = g$, $c_\ell = c$, $K_\ell = K$, $\rho_\ell = \rho$, $A_\ell = A$, so $\eta_\ell = \eta$ constant), Equation~\eqref{eq:recursive_precision} becomes $\tau^*_\ell = \rho \tau^*_{\ell-1} + C$ for some constant $C > 0$, which converges to $\bar{\tau} = C/(1-\rho) > 0$ as $\ell \to \infty$. In this case, the precision path converges to a \emph{positive limit}: adding more layers continues to increase precision (albeit with diminishing returns), and infinite depth is optimal whenever API costs are below a threshold. The boundedness of depth, if any, comes solely from cost considerations, not from fundamental information limits.

When $\eta_\ell < 1$, the dynamics change qualitatively. The following proposition characterizes the precision path under the information depreciation model. To maintain analytical transparency, we state the result for the uniform-technology case with potentially non-uniform attention budgets; the special case of uniform budgets ($A_\ell = A$, so $\eta_\ell = \eta$) recovers the baseline formula.

\begin{proposition}[Precision Path Characterization]\label{prop:precision_path}
Under the information depreciation model (Assumptions~\ref{ass:state_signals}--\ref{ass:endogenous_eta} with $\eta_\ell \in (\underline{\eta}, \bar{\eta})$) and uniform technology ($g_\ell = g$, $c_\ell = c$, $K_\ell = K$, $\rho_\ell = \rho$ for all $\ell$), define the composite signal contribution constant
\begin{equation}\label{eq:C_constant}
C \equiv \sum_{k=1}^{K} g\bigl(\alpha^{*}_{k}\bigr) \tau^{0}_{k},
\end{equation}
where $\alpha^{*}_{k}$ is the optimal attention allocation to source $k$. Let $\bar{\eta} \equiv \max_\ell \eta_\ell$ denote the maximal preservation rate across layers. The equilibrium posterior precision at layer $\ell$ satisfies
\begin{equation}\label{eq:recursive_tau_unified}
\tau^{*}_{\ell} = \rho \cdot \tau^{*}_{\ell-1} + C \eta_\ell^{\,\ell}, \qquad \tau^{*}_{0} = \sum_{k=1}^{K} \tau^{0}_{k}.
\end{equation}
Under uniform attention budgets ($A_\ell = A$, so $\eta_\ell = \eta$ for all $\ell$), the closed-form solution is
\begin{equation}\label{eq:tau_star_L}
\tau^{*}_{L} = \rho^{L} \tau^{*}_{0} + \frac{C\eta\bigl(\eta^{L} - \rho^{L}\bigr)}{\eta - \rho} \quad \text{for } \eta \neq \rho,
\end{equation}
and $\tau^{*}_{L} = \eta^{L}\bigl(\tau^{*}_{0} + CL\bigr)$ when $\eta = \rho$. Moreover:

\begin{enumerate}[label=(\roman*)]
\item \textbf{If $\eta < \rho$:} $\tau^{*}_{\ell}$ is hump-shaped. It increases for early layers, reaches a maximum at some finite $\ell^{*} \geq 1$, then decreases monotonically. The limiting precision is $\lim_{\ell \to \infty} \tau^{*}_{\ell} = 0$.

\item \textbf{If $\eta > \rho$:} $\tau^{*}_{\ell}$ is eventually decreasing. It may increase initially depending on initial conditions, but eventually declines monotonically to $0$.

\item \textbf{If $\eta = \rho$:} $\tau^{*}_{\ell} = \eta^{\ell}(\tau^{*}_{0} + C\ell)$, which is eventually decreasing to $0$.
\end{enumerate}
In all cases, $\lim_{\ell \to \infty} \tau^{*}_{\ell} = 0$.
\label{prop:geometric_decay}%
\end{proposition}

\begin{proof}
\noindent\textbf{Step 1: Derivation of the recurrence.}
Under uniform technology and Assumptions~\ref{ass:state_signals}--\ref{ass:endogenous_eta}, the new signal contribution at layer $\ell$ is $\sum_{k=1}^{K} g(\alpha^*_k) \tau^0_k \eta_\ell^{\,\ell} = C\eta_\ell^{\,\ell}$, where $C$ is defined in \eqref{eq:C_constant}. Substituting into \eqref{eq:recursive_precision} yields \eqref{eq:recursive_tau_unified}. Under uniform attention budgets ($A_\ell = A$), we have $\eta_\ell = \eta$ for all $\ell$, and the recurrence simplifies to $\tau^{*}_{\ell} = \rho \tau^{*}_{\ell-1} + C\eta^{\ell}$.

\noindent\textbf{Step 2: Closed-form solution (uniform budgets).}
Under uniform budgets, \eqref{eq:recursive_tau_unified} becomes a first-order linear non-homogeneous difference equation with constant coefficients.

\textit{Case 1: $\eta \neq \rho$.} The homogeneous solution is $\tau^{(h)}_{\ell} = A \rho^{\ell}$. Guessing a particular solution $\tau^{(p)}_{\ell} = B \eta^{\ell}$ and substituting:
\begin{equation*}
B\eta^{\ell} = \rho B \eta^{\ell-1} + C\eta^{\ell}
\quad\Longrightarrow\quad
B = \frac{C\eta}{\eta - \rho}.
\end{equation*}
Applying $\tau^{*}_{0} = A + B$ and solving for $\tau^{*}_{L}$ yields \eqref{eq:tau_star_L}.

\textit{Case 2: $\eta = \rho$.} The recurrence becomes $\tau^{*}_{\ell} = \eta \tau^{*}_{\ell-1} + C\eta^{\ell}$. Defining $u_{\ell} \equiv \tau^{*}_{\ell}/\eta^{\ell}$, we have $u_{\ell} = u_{\ell-1} + C$, which telescopes to $u_{\ell} = \tau^{*}_{0} + C\ell$. Hence $\tau^{*}_{L} = \eta^{L}(\tau^{*}_{0} + CL)$.

\noindent\textbf{Step 3: Vanishing precision (all cases).}
Since $\rho \in (0,1)$ and $\eta_\ell \leq \bar{\eta} < 1$ for all $\ell$ (by Assumption~\ref{ass:endogenous_eta}), both $\rho^{L} \to 0$ and $\eta_\ell^{\,L} \to 0$. In all three cases, $\lim_{L\to\infty} \tau^{*}_{L} = 0$. The critical difference from the $\eta = 1$ case is that $\bar{\eta} < 1$ is a \emph{structural bound}: even with unlimited attention, the preservation rate cannot reach unity, and geometric decay is inevitable.

\noindent\textbf{Step 4: Shape characterization.}
From \eqref{eq:recursive_tau_unified}, the layer-to-layer change is
\begin{equation}\label{eq:delta_tau}
\Delta_{\ell} \equiv \tau^{*}_{\ell} - \tau^{*}_{\ell-1} = -(1-\rho)\tau^{*}_{\ell-1} + C\eta_\ell^{\,\ell}.
\end{equation}
Under uniform budgets ($\eta_\ell = \eta$):

\textit{Subcase (i): $\eta < \rho$.} The $\rho^{\ell}$ term dominates the closed-form solution for large $\ell$. Since $\rho > \eta$, the negative term $-(1-\rho)\tau^{*}_{\ell-1}$ (which decays as $\rho^{\ell-1}$) eventually dominates the positive term $C\eta^{\ell}$ (which decays as $\eta^{\ell}$). Thus $\Delta_{\ell} < 0$ for all $\ell \geq \bar{L}_1$. For small $\ell$, however, the positive term may dominate if $\tau^{*}_{0}$ is small relative to $C\eta$, causing an initial increase. The path is therefore hump-shaped.

\textit{Subcase (ii): $\eta > \rho$.} As $\ell \to \infty$, $\tau^{*}_{\ell-1} \sim C\eta^{\ell}/(\eta-\rho)$. Substituting into \eqref{eq:delta_tau}:
\begin{equation*}
\Delta_{\ell} \sim C\eta^{\ell}\left(1 - \frac{1-\rho}{\eta-\rho}\right) = C\eta^{\ell}\frac{\eta - 1}{\eta - \rho} < 0,
\end{equation*}
since $\eta < 1$ and $\eta > \rho$. The sequence is eventually decreasing.

\textit{Subcase (iii): $\eta = \rho$.} Using $\tau^{*}_{\ell} = \eta^{\ell}(\tau^{*}_{0} + C\ell)$:
\begin{equation*}
\Delta_{\ell} = \eta^{\ell-1}\bigl[(\eta-1)\tau^{*}_{0} + C - C(1-\eta)\ell\bigr].
\end{equation*}
Since $\eta < 1$, the bracket is affine in $\ell$ with negative slope $-C(1-\eta) < 0$. Hence $\Delta_{\ell} < 0$ whenever $\ell > \tau^{*}_{0}/C + 1/(1-\eta)$.
\end{proof}

\begin{proposition}[Endogenous Preservation and the $\bar{\eta}$ Bound]\label{prop:eta_bar_bound}
Under Assumption~\ref{ass:endogenous_eta}, the preservation rate $\eta_\ell$ satisfies:
\begin{enumerate}[label=(\roman*)]
\item \textbf{Attention monotonicity:} $\eta_\ell$ is strictly increasing and concave in $A_\ell$:
\begin{equation}\label{eq:eta_increasing}
\frac{\partial \eta_\ell}{\partial A_\ell} = \frac{(\bar{\eta} - \underline{\eta}) \, a / K_\ell}{(\bar{\alpha}_\ell + a)^2} > 0, \qquad \frac{\partial^2 \eta_\ell}{\partial A_\ell^2} < 0.
\end{equation}

\item \textbf{Irreducible depreciation:} Even with unlimited attention, depreciation cannot be eliminated:
\begin{equation}\label{eq:eta_bar_limit}
\lim_{A_\ell \to \infty} \eta_\ell = \bar{\eta} < 1.
\end{equation}
Consequently, for any attention budget profile $\{A_\ell\}$, the cumulative precision at the top layer satisfies
\begin{equation}\label{eq:tau_L_upper_bound}
\tau^{*}_L \leq \rho^{L} \tau^{*}_{0} + \frac{C\bar{\eta}\bigl(\bar{\eta}^{L} - \rho^{L}\bigr)}{\bar{\eta} - \rho} \quad (\bar{\eta} \neq \rho),
\end{equation}
and the upper bound converges to $0$ as $L \to \infty$.

\item \textbf{Finite depth robustness:} The finite-optimal-depth result (Proposition~\ref{prop:optimal_depth}) holds \emph{regardless} of the attention budget profile. Even if $A_\ell \to \infty$ at every layer, the bound $\bar{\eta} < 1$ ensures $\lim_{L \to \infty} \tau^{*}_L = 0$.
\end{enumerate}
\end{proposition}

\begin{proof}
\noindent\textbf{Part (i):} Direct differentiation of \eqref{eq:endogenous_eta} with respect to $A_\ell$, using $\bar{\alpha}_\ell = A_\ell / K_\ell$, yields \eqref{eq:eta_increasing}.

\noindent\textbf{Part (ii):} Taking the limit of \eqref{eq:endogenous_eta} as $A_\ell \to \infty$:
\begin{equation*}
\lim_{A_\ell \to \infty} \eta_\ell = \underline{\eta} + (\bar{\eta} - \underline{\eta}) \cdot \lim_{\bar{\alpha}_\ell \to \infty} \frac{\bar{\alpha}_\ell}{\bar{\alpha}_\ell + a} = \underline{\eta} + (\bar{\eta} - \underline{\eta}) \cdot 1 = \bar{\eta}.
\end{equation*}
Since $\eta_\ell \leq \bar{\eta}$ for all $\ell$, the signal contribution at layer $\ell$ is bounded above by $C\bar{\eta}^{\,\ell}$. Substituting this upper bound into the recurrence \eqref{eq:recursive_tau_unified} yields \eqref{eq:tau_L_upper_bound}. Since $\bar{\eta} < 1$, the upper bound converges to $0$ as $L \to \infty$.

\noindent\textbf{Part (iii):} By Part (ii), $\tau^{*}_L$ is bounded above by a quantity that converges to $0$, so $\tau^{*}_L \to 0$ for \emph{any} budget profile. The proof of Proposition~\ref{prop:optimal_depth} then applies unchanged: there exists a finite $\bar{L}$ beyond which the marginal benefit of adding a layer is negative.
\end{proof}

\begin{remark}[Why $\bar{\eta} < 1$: Micro-foundations]\label{rem:eta_bar_microfoundation}
The bound $\bar{\eta} < 1$ is not an exogenous assumption; it reflects three fundamental limits of LLM information processing that no amount of attention can overcome:
\begin{enumerate}[label=(\roman*)]
\item \textbf{Irreversible embedding compression.} Each forward pass projects discrete tokens into a continuous embedding space through a non-invertible mapping. The entropy reduction is structural: a 4096-dimensional embedding of a 100,000-token vocabulary inherently discards information that cannot be recovered by downstream layers, regardless of context-window size.
\item \textbf{Non-invertible attention operations.} The attention mechanism computes weighted averages of value vectors. This is a many-to-one mapping: different input configurations can produce identical attention outputs. The operation discards positional and structural information about \emph{which} tokens contributed \emph{how much}, and this information loss is fundamental to the architecture.
\item \textbf{Lossy summarization as a design feature.} LLMs are trained to compress and generalize, not to preserve. The training objective (next-token prediction with KL regularization) explicitly rewards discarding idiosyncratic detail in favor of distributional patterns. Even with unlimited compute and context, the model's inductive bias favors abstraction over fidelity.
\end{enumerate}
These three forces explain why $\bar{\eta} < 1$ is a \emph{technological primitive} of transformer-based processing, not a modeling assumption that could be relaxed by better engineering. The bound $\bar{\eta}$ can be calibrated from empirical measurements of signal fidelity as a function of context length; see Section~\ref{sec:calibration}.
\end{remark}

\begin{remark}[What the Compounding Effect Is and Is Not]\label{rem:compounding_qualitative}
An earlier version of this paper claimed that $\tau^*_L$ is \emph{strictly decreasing} in $L$ for any fixed budget (a ``compounding effect''). That claim is incorrect. Proposition~\ref{prop:precision_path} shows that $\tau^*_\ell$ can \emph{increase} for early layers when fresh signals are abundant ($C\eta_\ell^{\,\ell}$ large relative to inherited precision). The correct statement is that precision \emph{eventually} declines due to information depreciation, not that it decreases monotonically from the first layer. The eventual decline to zero is the economically relevant property for establishing finite optimal depth.
\end{remark}

\subsection{Optimal Organizational Depth}
\label{subsec:optimal_depth}

Proposition~\ref{prop:precision_path} establishes that the precision path eventually declines to zero under information depreciation. We now show that this vanishing precision implies a finite optimal depth.

The principal chooses depth $L$ and the budget profile $\{A_\ell\}_{\ell=0}^{L-1}$ to maximize the value of the final decision. The optimization problem is:
\begin{equation}\label{eq:optimal_depth}
\max_{L, \{A_\ell\}} V_L(\tau^*_L) - \sum_{\ell=0}^{L-1} c_\ell(A_\ell).
\end{equation}

\begin{proposition}[Optimal Depth Existence]\label{prop:optimal_depth}
\label{prop:depth_finite}%
\label{prop:finite_depth}%
Under the information depreciation model (Assumptions~\ref{ass:hierarchy}--\ref{ass:endogenous_eta} with $\bar{\eta} < 1$), there exists a finite optimal depth $L^* < \infty$ for any attention budget profile $\{A_\ell\}$.
\end{proposition}

\begin{proof}
The proof has two parts. First, we show that $\tau^{*}_{\ell} \to 0$ as $\ell \to \infty$, which follows directly from Proposition~\ref{prop:precision_path} and the bound $\eta_\ell \leq \bar{\eta} < 1$ (Proposition~\ref{prop:eta_bar_bound}). Second, we show that this vanishing precision implies a finite optimal depth.

Since $V_L(\cdot)$ is strictly increasing and continuous, $\lim_{\ell \to \infty} V_L(\tau^{*}_{\ell}) = V_L(0)$. By Proposition~\ref{prop:precision_path}, there exists a finite threshold $\bar{L}$ such that $\tau^{*}_{\ell} < \tau^{*}_{\ell-1}$ for all $\ell \geq \bar{L}$. Since $V_L$ is strictly increasing, $V_L(\tau^{*}_{\ell}) < V_L(\tau^{*}_{\ell-1})$ for all $\ell \geq \bar{L}$.

For any $\ell \geq \bar{L}$, the net marginal benefit of adding layer $\ell$ is:
\begin{equation*}
\underbrace{V_L(\tau^{*}_{\ell}) - V_L(\tau^{*}_{\ell-1})}_{< 0} - \underbrace{c_{\ell}(A_{\ell})}_{\geq 0} < 0.
\end{equation*}
Thus no planner would ever choose $\ell > \bar{L}$, and $L^{*} \leq \bar{L} < \infty$ regardless of how small the marginal attention cost $\kappa$ is or how large the attention budgets $\{A_\ell\}$ are.

The critical mechanism is \emph{information depreciation bounded away from unity}: because $\eta_\ell \leq \bar{\eta} < 1$ for all $\ell$ (even with unlimited attention), each layer's signals carry less raw precision than the previous layer's ($\bar{\eta}^{\,\ell} \to 0$), while the attention budget can only raise $\eta_\ell$ toward $\bar{\eta}$ but never eliminate the geometric decay. The compounding of these two forces---inherent algorithmic noise (captured by $\bar{\eta} < 1$) and endogenous precision leakage ($\rho < 1$)---drives $\tau^{*}_{\ell}$ to zero and bounds the efficient depth of the chain.
\end{proof}

\vspace{1em}
\noindent\textbf{Necessity of the $\bar{\eta} < 1$ bound.}
The finite-depth result relies critically on Assumption~\ref{ass:endogenous_eta} with $\bar{\eta} < 1$. If $\bar{\eta} = 1$ (unbounded preservation with unlimited attention), then with sufficiently large $A_\ell$, $\eta_\ell$ could be driven arbitrarily close to $1$, and the recurrence \eqref{eq:recursive_tau_unified} would approach $\tau^{*}_{\ell} = \rho \tau^{*}_{\ell-1} + C$, with solution $\tau^{*}_{L} = \rho^{L}\tau^{*}_{0} + C(1-\rho^{L})/(1-\rho) \to C/(1-\rho) > 0$. The precision would converge to a positive limit, so $V(\tau^{*}_{L})$ converges to $V(C/(1-\rho)) > V(0)$ from below, and infinite depth could be optimal whenever the marginal cost $\kappa$ is below the limiting marginal benefit.

With the endogenous preservation rate bounded by $\bar{\eta} < 1$, this possibility is ruled out \emph{by construction}. Even as $A_\ell \to \infty$, we have $\eta_\ell \to \bar{\eta} < 1$, so the signal contribution at layer $\ell$ is bounded above by $C\bar{\eta}^{\,\ell}$, which decays geometrically. The bound $\bar{\eta} < 1$ is therefore \emph{sufficient} for the finite-depth result, \emph{regardless} of the attention cost or budget. This is the key sense in which endogenizing $\eta$ strengthens rather than weakens the model: the finite-depth result becomes robust to arbitrarily large attention budgets because inherent algorithmic noise places a technological ceiling on signal preservation.

\vspace{1em}
\noindent\textbf{Economic Discussion: Information Capital Depreciation.}

The critical insight of Proposition~\ref{prop:optimal_depth} is that \emph{depth-induced signal degradation} fundamentally alters the returns to layering. The endogenous preservation rate $\eta_\ell$ serves a role analogous to depreciation in capital accumulation: just as physical capital loses productive value through wear and tear, information capital (posterior precision) erodes through repeated algorithmic processing.

The preservation rate $\eta_\ell$ can be micro-founded through several channels:
\begin{enumerate}[label=(\roman*)]
\item \emph{Compressive loss:} Each LLM distills verbose inputs into concise outputs, inevitably discarding fine-grained detail. The summary statistics that propagate upward retain the gist but lose nuance.
\item \emph{Hallucination risk:} Models introduce plausible-sounding but factually incorrect content, injecting algorithmic noise that dilutes the signal-to-noise ratio.
\item \emph{Context window saturation:} Deeper layers must accommodate progressively more compressed representations; even with expanded context windows, the \emph{effective} precision per token declines.
\end{enumerate}

The endogenous framework clarifies the relationship between attention and depreciation. An $\eta_\ell$ close to $\bar{\eta}$ (achieved with large $A_\ell$) indicates high-fidelity transmission (e.g., retrieval-augmented systems with precise citations), while an $\eta_\ell$ close to $\underline{\eta}$ (small $A_\ell$) corresponds to lossy processing that obscures original signal content. Crucially, because $\bar{\eta} < 1$, even infinite attention cannot eliminate depreciation entirely. Our result that $L^{*}$ is bounded even when API calls are cheap ($\kappa \approx 0$) and attention budgets are large formalizes the intuition that organizational limits to AI delegation derive not from budget constraints alone, but from \emph{fundamental information-theoretic limits} inherent to transformer computation.

It is important to be precise about what drives the result. \emph{Attention scarcity alone} (finite $A_\ell$) is insufficient to guarantee finite optimal depth when costs vanish: if $\bar{\eta} = 1$, then with $A_\ell \to \infty$ we would have $\eta_\ell \to 1$, the precision would converge to a positive limit, and arbitrarily deep chains could be optimal. Conversely, \emph{information depreciation alone} ($\bar{\eta} < 1$) with unlimited attention budgets still bounds depth, as each layer's signals carry vanishingly little information no matter how much attention is applied. The two forces are \emph{complementary}: attention scarcity determines \emph{where} $\eta_\ell$ lies in the interval $[\underline{\eta}, \bar{\eta}]$, while the bound $\bar{\eta} < 1$ ensures that the \emph{total} information content always shrinks as the chain lengthens.

\subsection{Comparative Statics of Optimal Depth}
\label{subsec:comp_statics_depth}

We conclude with two comparative-statics results that link the optimal depth to economic primitives.

\begin{proposition}[Attention Budget and Depth]\label{prop:budget_depth}
The optimal depth $L^*$ is increasing in the attention budget profile:
\begin{equation}\label{eq:budget_depth}
\frac{\partial L^*}{\partial A_\ell} > 0 \quad \text{for all } \ell.
\end{equation}
\end{proposition}

\begin{proof}
A larger attention budget at any layer raises the transmission factor $\rho_\ell$, reducing the precision loss per layer. This makes deeper chains more attractive because the marginal benefit of adding a layer falls less quickly relative to the marginal cost. Formally, by Proposition~\ref{prop:precision_path}, a higher $\rho$ (from larger $A_\ell$) raises the precision path $\tau^*_\ell$ at all layers and can delay the threshold $\bar{L}$ beyond which precision declines, allowing a larger optimal $L^*$.
\end{proof}

\begin{proposition}[Cost of Attention and Depth]\label{prop:cost_depth}
The optimal depth $L^*$ is decreasing in the marginal cost of attention:
\begin{equation}\label{eq:cost_depth}
\frac{\partial L^*}{\partial \kappa_\ell} < 0 \quad \text{for all } \ell,
\end{equation}
where $\kappa_\ell$ parameterizes the cost function $c_\ell(A_\ell) = \kappa_\ell A_\ell + \frac{\phi_\ell}{2} A_\ell^2$.
\end{proposition}

\begin{proof}
Higher API costs make each layer more expensive, pushing the organization toward shallower structures. This is the attention-allocation analog of the Coasian prediction that higher coordination costs favor integration.
\end{proof}

\vspace{1em}
\noindent\textbf{Summary.} Proposition~\ref{prop:front_loading} establishes that the optimal architecture front-loads attention capacity: larger context windows should be deployed at early layers. Proposition~\ref{prop:optimal_depth} shows that finite optimal depth is a generic property of the information depreciation model. Propositions~\ref{prop:budget_depth} and~\ref{prop:cost_depth} provide the comparative-statics predictions: larger context windows permit deeper chains, while higher API costs push toward shallower ones. These four results collectively discipline the design of multi-agent AI architectures and generate testable predictions about the relationship between technology parameters and organizational form.
